% 《金融研究》(Journal of Financial Research) LaTeX 模板
% 适配 GB/T 7714-2015 参考文献格式
% 格式：A4，双栏，中文，摘要+关键词
%
% 编译方式：xelatex -> bibtex -> xelatex -> xelatex

% ─── 文档类 ────────────────────────────────────────────────────────────────
\documentclass[10pt, UTF8, twocolumn]{ctexart}
% 10pt: 金融研究正文字号
% twocolumn: 双栏排版

% ─── 页面布局 ─────────────────────────────────────────────────────────────
\usepackage[
    a4paper,
    top=22mm,
    bottom=22mm,
    left=20mm,
    right=20mm,
    headheight=10mm,
    footskip=8mm,
    columnsep=8mm,
]{geometry}

% ─── 语言与字体 ────────────────────────────────────────────────────────────
\usepackage{xeCJK}
%\setCJKmainfont{...}[...]  % 系统字体配置（可选，ctex 自动 fallback）

% ─── 数学与符号 ────────────────────────────────────────────────────────────
\usepackage{amsmath, amssymb}
\usepackage{mathtools}
\usepackage{booktabs}
\usepackage{longtable}
\usepackage{graphicx}
\graphicspath{{./figures/}}

% ─── 间距控制 ─────────────────────────────────────────────────────────────
\usepackage{setspace}
%\setstretch{1.5}
\setlength{\parindent}{2em}
\setlength{\parskip}{0pt}

% ─── 页面样式 ─────────────────────────────────────────────────────────────
\usepackage{fancyhdr}
\fancypagestyle{JRYJ}{
    \fancyhf{}
    \fancyhead[C]{\thepage}
    \fancyfoot[C]{\thepage}
    \fancyhead[LO,RE]{}
    \fancyhead[CO,CE]{}
    \renewcommand{\headrulewidth}{0pt}
    \renewcommand{\footrulewidth}{0pt}
}
\pagestyle{JRYJ}

% ─── 超链接 ───────────────────────────────────────────────────────────────
\usepackage[colorlinks=true, linkcolor=black, citecolor=black, urlcolor=black]{hyperref}

% ─── 参考文献（GB/T 7714-2015 numeric）─────────────────────────────────────
\usepackage[numbers,sort&compress]{natbib}
\bibliographystyle{gbt7714-plain}
% 如有 gbt7714-2005.bst 可改用: \bibliographystyle{gbt7714-2005}

% ─── 标题宏 ───────────────────────────────────────────────────────────────
\makeatletter
\def\@JRYJtitle#1{\gdef\@JRYJ@title{#1}}\def\@JRYJ@title{}
\def\@JRYJauthor#1{\gdef\@JRYJ@author{#1}}\def\@JRYJ@author{}
\def\@JRYJaffiliation#1{\gdef\@JRYJ@affiliation{#1}}\def\@JRYJ@affiliation{}
\def\@JRYJabstract#1{\gdef\@JRYJ@abstract{#1}}\def\@JRYJ@abstract{}
\def\@JRYJkeywords#1{\gdef\@JRYJ@keywords{#1}}\def\@JRYJ@keywords{}
\makeatother

% 摘要环境
\NewDocumentEnvironment{jryjabstract}{}{
    \par\noindent\textbf{摘\quad 要}\xiaosi
    \wuhao
}{
    \par\vspace{\baselineskip}
}

% 关键词环境
\NewDocumentEnvironment{jryjkeywords}{}{
    \par\noindent\textbf{关键词}\xiaosi\ \wuhao
}{
    \par\vspace{0.5\baselineskip}
}

% 设置命令
\NewDocumentCommand{\JRYJsettitle}{m}{\@JRYJtitle{#1}}
\NewDocumentCommand{\JRYJsetauthor}{m}{\@JRYJauthor{#1}}
\NewDocumentCommand{\JRYJsetaffiliation}{m}{\@JRYJaffiliation{#1}}
\NewDocumentCommand{\JRYJsetabstract}{m}{\@JRYJabstract{#1}}
\NewDocumentCommand{\JRYJsetkeywords}{m}{\@JRYJkeywords{#1}}

% 字号命令
\newcommand{\xiaosi}{\fontsize{12pt}{18pt}\selectfont}
\newcommand{\wuhao}{\fontsize{9pt}{13pt}\selectfont}
% 标题字号
\newcommand{\JRYJtitlefont}{\fontsize{18pt}{26pt}\bfseries\selectfont}
% 作者字号
\newcommand{\JRYJauthorfont}{\fontsize{10.5pt}{16pt}\selectfont}

% ─── 正文格式 ─────────────────────────────────────────────────────────────
% 一级标题：黑体加粗，居中
\ctexset{section/format={\centering\zihao{-3}\CJKfamily{zhhei}\bfseries}}
% 二级标题：黑体
\ctexset{subsection/format={\zihao{4}\CJKfamily{zhhei}}}
% 三级标题：黑体
\ctexset{subsubsection/format={\zihao{-4}\CJKfamily{zhhei}}}

% ─── 图表浮动 ─────────────────────────────────────────────────────────────
\usepackage{caption}
\captionsetup{font=small, skip=6pt}

% ─── 定理环境 ─────────────────────────────────────────────────────────────
\newtheorem{hypothesis}{假设}
\newtheorem{assumption}{假定}
\newtheorem{finance}{命题}
\newtheorem*{proof}{证明}

% ─────────────────────────────────────────────────────────────────────────
\begin{document}

% ══════════════════════════════════════════════════════════════════════════════
% 封面信息（盲审时删除作者和单位信息）
% ══════════════════════════════════════════════════════════════════════════════

\JRYJsettitle{论文标题（中文，不超过20字）}

% 盲审时注释掉以下两行
\JRYJsetauthor{作者一$^{1}$~~作者二$^{2}$~~作者三$^{1}$}
\JRYJsetaffiliation{$^{1}$北京大学光华管理学院~~$^{2}$清华大学五道口金融学院}

\begin{jryjabstract}
本文基于XX数据，采用XX方法，实证研究了XX问题。研究发现：（1）……；（2）……；（3）……。本文的研究对理解……具有理论意义，对金融监管和金融市场参与者……具有实践价值。
\end{jryjabstract}

\begin{jryjkeywords}
关键词一；关键词二；关键词三；关键词四
\end{jryjkeywords}

% ══════════════════════════════════════════════════════════════════════════════
% 正文
% ══════════════════════════════════════════════════════════════════════════════

\section{引言}
\label{sec:intro}

本文的研究背景与问题、意义和主要贡献如下。

第一，……

第二，……

第三，……

\section{文献综述}
\label{sec:lit}

\subsection{理论文献回顾}

\subsection{实证文献回顾}
已有文献在……方面做出了重要贡献，但仍存在以下不足：（1）……；（2）……；（3）……。

本文与已有研究的关键区别在于：……。

\section{研究设计}
\label{sec:design}

\subsection{样本与数据}
本文以……为研究样本，时间跨度为……至……。数据来源包括……数据库、……数据库和……。

样本筛选过程：（1）剔除……；（2）剔除……；（3）剔除……。

\subsection{变量定义}

\subsubsection{被解释变量}
$Y_{it}$ 表示……，数据来源于……。

\subsubsection{解释变量}
核心解释变量$X_{it}$ 为……。该指标由……构造得到。

工具变量$IV_{it}$ 为……。选择依据：……。

\subsubsection{控制变量}
参考已有研究（\cite{xxx}），本文还控制了以下变量：$Z_1$（……）、$Z_2$（……）、年份固定效应等。

\begin{table}[htbp]
\centering
\caption{变量定义表}
\label{tab:var}
\xiaosi
\begin{tabular}{p{2cm}p{2.5cm}p{5.5cm}}
\toprule
变量类型 & 变量名称 & 变量定义 \\
\midrule
因变量 & $Y$ & ……（数据来源：XX数据库）\\
\midrule
自变量 & $X$ & ……（指标构造方法）\\
\midrule
控制变量 & $Z_1$ & ……（定义）\\
& $Z_2$ & ……（定义）\\
& $Year$ & 年份虚拟变量\\
\midrule
工具变量 & $IV$ & ……（选择依据）\\
\bottomrule
\end{tabular}
\end{table}

\subsection{模型构建}
为检验假设1，本文构建如下双向固定效应回归模型：

\begin{equation}
\label{eq:baseline}
Y_{it} = \alpha + \beta X_{it} + \gamma Z_{it} + \mu_i + \lambda_t + \varepsilon_{it}
\end{equation}

其中，$X_{it}$ 为核心解释变量，$Z_{it}$ 为控制变量向量，$\mu_i$ 和 $\lambda_t$ 分别为企业和年份固定效应，$\varepsilon_{it}$ 为随机误差项。

\section{实证结果}
\label{sec:results}

\subsection{描述性统计}
表~\ref{tab:desc} 报告了主要变量的描述性统计。

\begin{longtable}{ccccccc}
\caption{描述性统计}
\label{tab:desc}
\\\toprule
变量 & 观测数 & 均值 & 标准差 & 最小值 & 中位数 & 最大值 \\
\midrule
\endfirsthead
\multicolumn{7}{c}{\textbf{续表\thetable\ 描述性统计}}\\
\toprule
变量 & 观测数 & 均值 & 标准差 & 最小值 & 中位数 & 最大值 \\
\midrule
\endhead
\bottomrule
\endfoot
$Y$ & 1000 & 0.05 & 0.20 & -0.50 & 0.03 & 0.60 \\
$X$ & 1000 & 0.30 & 0.15 & 0.10 & 0.28 & 0.70 \\
$Z_1$ & 1000 & 5.20 & 1.50 & 1.00 & 5.00 & 9.00 \\
\bottomrule
\end{longtable}

\subsection{基准回归分析}
表~\ref{tab:reg} 报告了基准回归结果。

\begin{table}[htbp]
\centering
\caption{基准回归结果}
\label{tab:reg}
\xiaosi
\begin{tabular}{lccc}
\toprule
 & (1) & (2) & (3) \\
\midrule
$X$ & 0.05** & 0.06*** & 0.04** \\
    & (0.02) & (0.02) & (0.02) \\
$Z_1$ & & 0.10*** & 0.08** \\
      &      & (0.03)  & (0.03) \\
\midrule
常数项 & 0.01 & 0.02 & 0.03 \\
       & (0.02) & (0.02) & (0.02) \\
\midrule
控制变量 & 否 & 是 & 是 \\
企业固定效应 & 否 & 否 & 是 \\
年份固定效应 & 否 & 否 & 是 \\
\midrule
观测数 & 1000 & 1000 & 1000 \\
$R^2$ & 0.15 & 0.18 & 0.22 \\
\bottomrule
\multicolumn{4}{l}{\small 注：*** p<0.01, ** p<0.05, * p<0.1。括号内为聚类稳健标准误（聚类到企业层面）。}
\end{tabular}
\end{table}

\section{内生性处理}
\label{sec:endo}

\subsection{工具变量法}
为缓解反向因果导致的内生性问题，本文采用工具变量方法。

第一阶段：$X_{it} = \pi_0 + \pi_1 IV_{it} + \cdots + \varepsilon_{it}$

第二阶段：$Y_{it} = \alpha + \beta \hat{X}_{it} + \gamma Z_{it} + \mu_i + \lambda_t + \varepsilon_{it}$

工具变量有效性检验：弱工具变量检验F统计量为……，Stock-Yogo临界值为……，拒绝弱工具变量假设。

\subsection{其他内生性处理策略}
（1）滞后解释变量；（2）PSM倾向得分匹配；（3）Heckman两阶段模型。

\section{稳健性检验}
\label{sec:robust}

（1）替换核心变量：采用……替换$X$，结果依然显著。

（2）改变样本范围：剔除……样本后，主要结论不变。

（3）改变控制变量：加入/剔除……变量，结论稳健。

（4）改变固定效应结构：加入……交互固定效应。

\section{进一步分析}
\label{sec:further}

\section{结论}
\label{sec:concl}

本文基于……数据，实证研究了……问题。研究发现：（1）……；（2）……；（3）……。

理论贡献：第一，……；第二，……。

实践启示：对于金融监管机构，……；对于金融机构，……；对于投资者，……。

本文的局限性及未来研究方向：……。

\section*{参考文献}
\addcontentsline{toc}{section}{参考文献}
\bibliography{references}

% ══════════════════════════════════════════════════════════════════════════════
% 附录
% ══════════════════════════════════════════════════════════════════════════════
\clearpage
\appendix
\section*{附录}
\addcontentsline{toc}{section}{附录}

\end{document}
